\documentclass[12pt,a4paper]{article}

\PassOptionsToPackage{hyphens}{url}   % let long URLs break at / and -
\usepackage[a4paper,margin=1in]{geometry}
\usepackage{graphicx}
\usepackage{booktabs}
\usepackage{amsmath,amssymb}
\usepackage{float}
\usepackage{hyperref}
\usepackage{xcolor}
\usepackage{listings}
\usepackage{enumitem}
\usepackage{fancyhdr}
\usepackage{longtable}
\usepackage{array}

% Scale a table down to the text width only if it would otherwise overflow.
\newsavebox{\fitbox}
\newenvironment{fitwidth}
  {\begin{lrbox}{\fitbox}}
  {\end{lrbox}\ifdim\wd\fitbox>\linewidth\resizebox{\linewidth}{!}{\usebox{\fitbox}}\else\usebox{\fitbox}\fi}
% Let long code identifiers (e.g. load\_breast\_cancer) wrap at underscores instead of running into the margin.
\DeclareRobustCommand{\_}{\textunderscore\allowbreak}
\setlength{\emergencystretch}{3em}

\hypersetup{colorlinks=true,linkcolor=blue,urlcolor=blue}

\pagestyle{fancy}
\fancyhf{}
\lhead{ICS1512}
\rhead{Machine Learning Algorithms Laboratory}
\cfoot{\thepage}

\lstset{
language=Python,
basicstyle=\ttfamily\small,
numbers=left,
frame=single,
breaklines=true,
showstringspaces=false,
keywordstyle=\color{blue},
commentstyle=\color{gray},
stringstyle=\color{purple}
}

\begin{document}

\noindent\begin{tabular}{|p{4cm}|p{\dimexpr\linewidth-4cm-4\tabcolsep-3\arrayrulewidth\relax}|}
\hline
Experiment No. & 9 \\ \hline
Experiment Title & Perceptron vs Multilayer Perceptron (A/B Experiment) with Hyperparameter Tuning
\footnote{\textbf{GitHub Repository: }\url{https://github.com/Nightingales21/ICS1512-Machine-Learning-Labratory}}\\ \hline
Student Name & Nitin \\ \hline
Register Number & 312224700142 \\ \hline
Faculty & Dr. Poreddy Ajay Kumar Reddy \\ \hline
Submission Date & \today \\ \hline
\end{tabular}

\section{Aim and Objective}
To implement and compare two neural models on the same 62-class handwritten
character recognition task:
\begin{itemize}[nosep]
\item \textbf{Model A:} the Single-Layer Perceptron Learning Algorithm (PLA),
written from scratch with a step activation and the classical
$w \leftarrow w + \eta(y - \hat{y})x$ update rule;
\item \textbf{Model B:} a Multilayer Perceptron (MLP) with hidden layers and
nonlinear activations, trained by backpropagation.
\end{itemize}
Hyperparameters for Model B (activation function, optimizer, learning rate,
batch size, hidden-layer architecture and regularisation) are selected by
systematic tuning rather than by assertion, and the two models are then
compared as a controlled A/B experiment on an identical train--test split.

\section{Dataset Description}
\begin{longtable}{|p{6cm}|p{8cm}|}
\hline
Dataset Name & English Handwritten Characters \\ \hline
Dataset Source & Kaggle (\texttt{dhruvildave/\allowbreak english-\allowbreak handwritten-\allowbreak characters-\allowbreak dataset}),
supplied as \texttt{english.csv} plus an \texttt{Img/} folder of PNG files \\ \hline
Number of Samples & 3{,}410 images \\ \hline
Number of Classes & 62 (digits 0--9, uppercase A--Z, lowercase a--z) \\ \hline
Samples per Class & 55, exactly balanced \\ \hline
Number of Features & 1{,}024 after preprocessing ($32 \times 32$ grayscale pixels) \\ \hline
Missing Values & 0 (verified programmatically) \\ \hline
Train--Test Split & 80\% / 20\%, stratified, \texttt{random\_state=42}
(2{,}728 train / 682 test) \\ \hline
\end{longtable}

The dataset is perfectly balanced: every one of the 62 classes contributes
exactly 55 images, so accuracy is a meaningful headline metric and no class
weighting is required. The test split therefore contains only 11 images per
class, which is worth keeping in mind when reading per-class scores: a single
extra correct prediction moves a class's recall by 9 percentage points.

\textbf{Why this dataset is hard.} Unlike MNIST-style digit sets, the label
space here mixes digits with both letter cases, and several classes are
\emph{visually identical} once the character is isolated from its context:
\texttt{0}/\texttt{O}/\texttt{o}, \texttt{1}/\texttt{l}/\texttt{I},
\texttt{S}/\texttt{s}, \texttt{X}/\texttt{x} and \texttt{Z}/\texttt{z} differ
mainly in size, which is destroyed by resizing every image to a common
$32\times32$ grid. This places a hard ceiling on achievable accuracy that is
well below the 90\%+ typical of 10-class digit recognition, and it is the
single most important factor behind the absolute numbers reported below.

\section{Preprocessing Steps}
\begin{itemize}[nosep]
\item \textbf{Grayscale conversion:} each PNG is converted to 8-bit
grayscale (\texttt{PIL.\allowbreak Image.\allowbreak convert("L")}), discarding
colour, which carries no information about character identity.
\item \textbf{Resize:} images arrive at 1200$\times$900 and are resized to a
fixed $32\times32$ grid with bilinear interpolation, giving every sample the
same dimensionality.
\item \textbf{Normalisation:} pixel values are scaled to $[0,1]$ and then
\emph{inverted}, so ink $=1.0$ and paper $=0.0$. This matters for the
perceptron specifically: with the raw polarity the background (the large
majority of pixels) would carry the large values and dominate every dot
product.
\item \textbf{Flattening:} the $32\times32$ grid is flattened to a
1{,}024-dimensional vector, which is what both models consume.
\item \textbf{No missing-value handling or categorical encoding} was required;
labels were mapped to integer class indices $0..61$ in sorted order.
\end{itemize}

\noindent For reference, PCA on the training split needs \textbf{92 of the
1{,}024 components} to retain 95\% of pixel variance, confirming that the raw
pixel representation is highly redundant.

\section{Exploratory Data Analysis}
The reusable plotting helpers from \texttt{ml\_lab\_utils.py} (Experiment 1)
were used for the mandatory figure styling; the EDA panel below is specific to
image data rather than the tabular 12-panel summary used in Experiments 2--7.

\begin{figure}[H]
\centering
\includegraphics[width=\linewidth]{figures/eda_handwritten_characters.png}
\caption{EDA summary: sample characters, class distribution, mean image, ink
coverage, per-pixel variance, and a PCA projection coloured by character group.}
\label{fig:eda}
\end{figure}
\textbf{Inference:} The class distribution (panel 2) is exactly flat at 55
images per class. The mean image (panel 3) is a diffuse central blob, and the
per-pixel variance map (panel 5) shows that essentially all discriminative
signal sits in the central region -- the border pixels are near-constant and
contribute almost nothing, which is why 92 components already capture 95\% of
the variance. Most importantly, the PCA projection (panel 6) shows digits,
uppercase and lowercase letters occupying \emph{the same} region of feature
space with no visible separation: the classes are not linearly separable in
pixel space, which directly predicts the PLA's difficulty in
Section~\ref{sec:pla}.

\section{Experimental Setup}
\begin{center}
\begin{minipage}{\linewidth}\centering
\begin{fitwidth}
\begin{tabular}{ll}
\toprule
Python Version & 3.11.15 \\
Scikit-Learn Version & 1.9.0 \\
NumPy / Pandas Version & 2.4.6 / 3.0.5 \\
Pillow Version & 12.2.0 \\
IDE & Jupyter Notebook / VS Code \\
Operating System & Windows 11 \\
Hardware & Standard lab workstation (CPU only) \\
\bottomrule
\end{tabular}
\end{fitwidth}
\end{minipage}
\end{center}

\vspace{0.3cm}
\textbf{Environment activation command (as mandated by the lab manual):}
\begin{lstlisting}[language=bash]
source /opt/anaconda3/bin/activate anaconda-navigation
\end{lstlisting}

\section{Model A: Perceptron Learning Algorithm (from scratch)}
\label{sec:pla}

The PLA was implemented directly from the update rule in the lab manual, with
one binary perceptron per class (one-vs-rest) and a hard step activation.
Prediction takes the arg-max over the 62 perceptron scores.

\begin{lstlisting}[language=Python, caption={Core of the from-scratch PLA, from \texttt{experiment9\_pla\_vs\_mlp.py}.}]
for _ in range(self.n_epochs):
    order = rng.permutation(len(X))
    for i in order:                      # online (sample-by-sample) updates
        xi = X[i]
        pred = np.where(xi @ self.W + self.b >= 0.0, 1.0, -1.0)  # step activation
        err = Y[i] - pred                # 0 where correct, +-2 otherwise
        if np.any(err):
            self.W += self.eta * np.outer(xi, err)
            self.b += self.eta * err
\end{lstlisting}

\noindent\textbf{Note on initialisation.} Weights are initialised from a small
random normal ($\sigma = 0.01$) rather than zeros. With all-zero
initialisation, a one-vs-rest arg-max perceptron is exactly invariant to
$\eta$ (every weight is a sum of $\eta$-scaled updates, so $\eta$ cancels in
the arg-max), and a learning-rate study would be meaningless.

\begin{center}
\begin{minipage}{\linewidth}\centering
\textbf{Table 1: PLA Learning-Rate Sweep (40 epochs each)}\par\smallskip
\begin{fitwidth}
\begin{tabular}{lccc}
\toprule
Learning rate $\eta$ & Final Training Error & Test Accuracy & Time (s) \\
\midrule
0.001 & 0.3853 & 0.2229 & 23.4 \\
0.01  & 0.3849 & 0.2009 & 20.9 \\
\textbf{0.1} & \textbf{0.3728} & \textbf{0.2507} & 12.5 \\
1.0   & 0.4219 & 0.2067 & 12.3 \\
\bottomrule
\end{tabular}
\end{fitwidth}
\end{minipage}
\end{center}

\textbf{Inference:} $\eta = 0.1$ is the best of the four, reaching 25.07\%
test accuracy. The differences are modest (20--25\%) because the learning rate
is not the binding constraint here -- the model class is. Note the
non-monotonic pattern: $\eta = 1.0$ is \emph{worse} than $\eta = 0.1$, since
very large updates cause each correction to overshoot and undo previously
learned structure.

\begin{figure}[H]
\centering
\includegraphics[width=0.8\linewidth]{figures/pla_convergence.png}
\caption{PLA training error vs epochs for four learning rates.}
\label{fig:pla_conv}
\end{figure}
\textbf{Inference:} None of the four curves converges to zero training error.
They fall quickly for roughly ten epochs and then oscillate in a band around
0.37--0.42 without ever settling. This is the textbook signature of the
perceptron convergence theorem's precondition being violated: the theorem
guarantees termination \emph{only} for linearly separable data, and this data
(as panel 6 of Figure~\ref{fig:eda} showed) is not. The oscillation is the
algorithm perpetually fixing one class's mistakes at the expense of another's.

\section{Model B: MLP Hyperparameter Tuning}

Tuning was done in four sequential stages with \texttt{GridSearchCV} (3-fold,
accuracy). Each stage fixes the winner of the previous stage, which keeps the
search to 28 configurations instead of the 360 a full factorial grid would
need. Stage 4 was added \emph{after} observing that stages 1--3 had driven
training accuracy to almost 100\% while test accuracy stalled.

\begin{center}
\begin{minipage}{\linewidth}\centering
\textbf{Table 2: Staged Hyperparameter Search (3-fold CV accuracy)}\par\smallskip
\begin{fitwidth}
\begin{tabular}{llcl}
\toprule
Stage & Parameters searched & Best CV Acc. & Winner \\
\midrule
1 & activation $\times$ optimizer & 0.4278 & \texttt{relu}, \texttt{adam} \\
2 & learning rate $\times$ batch size & 0.4644 & \texttt{lr=0.01}, \texttt{batch=128} \\
3 & hidden-layer architecture & 0.4732 & \texttt{(512,)} \\
4 & \texttt{alpha} $\times$ \texttt{early\_stopping} & 0.4732 & \texttt{alpha=1e-4}, no early stop \\
\bottomrule
\end{tabular}
\end{fitwidth}
\end{minipage}
\end{center}

\begin{figure}[H]
\centering
\includegraphics[width=\linewidth]{figures/mlp_tuning_stages.png}
\caption{Cross-validated accuracy for every configuration tried in the four tuning stages.}
\label{fig:tuning}
\end{figure}
\textbf{Inference:} Stage 1 is the most decisive: \texttt{relu}+\texttt{adam}
beats every \texttt{sgd} variant by a wide margin, because plain SGD at a
fixed learning rate makes very slow progress on 1{,}024 sparse pixel inputs
where Adam's per-parameter scaling adapts automatically. Stage 2 adds the
next-largest gain (+3.7 points) -- the default $10^{-3}$ learning rate is too
small for this problem. Stage 3 shows width helps only marginally, and deeper
two-layer networks do not beat a single wide layer on a dataset of just 2{,}728
training images.

\begin{figure}[H]
\centering
\includegraphics[width=0.8\linewidth]{figures/mlp_loss_curve.png}
\caption{MLP training loss vs epochs, tuned configuration vs untuned baseline.}
\label{fig:mlp_loss}
\end{figure}

\section{A/B Comparison: PLA vs Tuned MLP}

\begin{center}
\begin{minipage}{\linewidth}\centering
\textbf{Table 3: A/B Comparison on the Held-Out Test Set (682 images, macro-averaged)}\par\smallskip
\begin{fitwidth}
\begin{tabular}{lccccc}
\toprule
Model & Accuracy & Precision & Recall & F1-score & Train Time (s) \\
\midrule
PLA (Model A) & 0.2507 & 0.2973 & 0.2507 & 0.2336 & 11.5 \\
MLP (Model B, tuned) & \textbf{0.4648} & \textbf{0.4904} & \textbf{0.4648} & \textbf{0.4647} & 13.1 \\
\bottomrule
\end{tabular}
\end{fitwidth}
\end{minipage}
\end{center}

\begin{figure}[H]
\centering
\includegraphics[width=\linewidth]{figures/model_comparison.png}
\caption{Left: headline metrics, PLA vs tuned MLP. Right: per-class F1-score across all 62 classes.}
\label{fig:comparison}
\end{figure}

\begin{figure}[H]
\centering
\includegraphics[width=\linewidth]{figures/roc_curves.png}
\caption{ROC curves (one-vs-rest), micro- and macro-averaged over the 62 classes.}
\label{fig:roc}
\end{figure}

\begin{figure}[H]
\centering
\includegraphics[width=\linewidth]{figures/confusion_matrices.png}
\caption{Confusion matrices ($62 \times 62$) for the PLA and the tuned MLP.}
\label{fig:cm}
\end{figure}

\section{Observations}
\begin{itemize}
\item \textbf{Why does PLA underperform compared to MLP?} Because the classes are not linearly
separable in pixel space. The PLA can only carve the 1{,}024-dimensional input
with 62 hyperplanes, one per class, and Figure~\ref{fig:eda} (panel 6) shows
the three character groups occupying one overlapping cloud with no linear
boundary between them. The MLP's hidden layer of 512 ReLU units first
\emph{re-represents} each image as 512 learned nonlinear features, and only
then applies a linear separation in that learned space -- which is exactly the
capability the single-layer model lacks. The gap is large and consistent
across every metric: accuracy 0.2507 vs 0.4648 (1.85$\times$), macro F1
0.2336 vs 0.4647 (1.99$\times$), and macro ROC-AUC 0.8216 vs 0.9259. A
secondary reason is the step activation itself: it has zero gradient
everywhere, so the PLA cannot express "how wrong" a prediction was, only that
it was wrong.
\item \textbf{Which hyperparameters had the most impact on MLP performance?}
The \textbf{learning rate} mattered most.
Stage 2 lifted cross-validated accuracy from 0.4278 to 0.4644 (+3.7 points),
almost all of it from raising the learning rate from the default $10^{-3}$ to
$10^{-2}$. \textbf{Activation and optimizer} (stage 1) were the next most
important -- every \texttt{sgd} configuration trailed its \texttt{adam}
counterpart badly. \textbf{Architecture} (stage 3) contributed under one point
(0.4644 $\rightarrow$ 0.4732), and \textbf{regularisation} (stage 4)
contributed nothing measurable at all.
\item \textbf{Did optimizer choice (SGD vs Adam) affect convergence?}
Yes, decisively. With everything else held
fixed, \texttt{adam} configurations reached roughly 0.43 cross-validated
accuracy while \texttt{sgd} configurations lagged well behind in the same
60-epoch budget. Adam maintains a per-parameter adaptive step size, which suits
1{,}024 pixel inputs whose individual variances differ by orders of magnitude
(the border pixels are near-constant, the central ones highly variable, as the
variance map in Figure~\ref{fig:eda} shows). Plain SGD applies one global step
size to all of them, so it is simultaneously too large for the high-variance
central pixels and too small for the rest.
\item \textbf{Did adding more hidden layers always improve results?}
No. Stage 3 tested two-layer networks
\texttt{(256,128)} and \texttt{(512,256)} against single hidden layers, and a
single wide layer \texttt{(512,)} won (0.4732). With only 2{,}728 training
images and 1{,}024 inputs, a second hidden layer adds parameters to a model
that is already capable of memorising the training set (it reaches 99.89\%
training accuracy) without adding usable generalisation. Depth pays off when
there is enough data to constrain it; here there is not.
\item \textbf{Did MLP show overfitting? How could it be mitigated?}
Yes, severely, and this is the single most
important caveat on the result. The tuned MLP reaches \textbf{99.89\%
training accuracy} but only \textbf{46.48\%} on the test set -- a
generalisation gap of over 53 points -- with the training loss driven down to
0.025. The honest evidence that tuning did not fix this is that the
\emph{untuned baseline} MLP (\texttt{lr}=$10^{-3}$, 256 units) scored
\textbf{48.83\%} on the test set, slightly \emph{better} than the tuned model,
even though it scored worse in cross-validation. In other words the staged
search bought a 4.5-point cross-validation gain that did not transfer, which
is what a 3-fold CV over 2{,}728 images with 62 classes (roughly 14 validation
images per class per fold) is prone to do.

Stage 4 attempted the two standard mitigations directly. Neither helped on
cross-validated accuracy: the L2 penalty was flat up to \texttt{alpha}=0.1
(0.4725 vs 0.4732) and clearly harmful at \texttt{alpha}=1.0 (0.4142), while
\texttt{early\_stopping=True} lost 3--5 points in every pairing because it
spends 10\% of an already small training set on an internal validation split.
The mitigations that would actually be expected to work here address the data
rather than the penalty: \textbf{data augmentation} (small rotations, shifts
and elastic distortions, which is standard for handwriting), \textbf{more
training data per class} (55 images per class is very few for 62 classes), or
a \textbf{convolutional architecture} whose weight sharing and translation
invariance encode the structure of images directly instead of forcing a dense
layer to learn it from scratch.
\end{itemize}

\section{Conclusion}
This experiment compared a from-scratch single-layer
Perceptron against a systematically tuned MLP on 62-class handwritten
character recognition, under an identical stratified 80/20 split.

The A/B comparison is unambiguous in direction: the MLP roughly doubles the
PLA on every metric (accuracy 0.4648 vs 0.2507, macro F1 0.4647 vs 0.2336,
macro ROC-AUC 0.9259 vs 0.8216) for essentially the same training cost (13.1\,s
vs 15.1\,s). The PLA's training-error curves never converge, which is the
perceptron convergence theorem behaving exactly as advertised on data that is
not linearly separable, and the PCA projection confirms that this data is not.
A hidden layer with a nonlinear activation is therefore not an optional
refinement on this task -- it is what makes the task learnable at all.

Two honest qualifications belong with that headline. First, the absolute
accuracies are low because the label space is intrinsically ambiguous: with
\texttt{0}/\texttt{O}/\texttt{o}, \texttt{1}/\texttt{l}/\texttt{I} and
\texttt{S}/\texttt{s} rendered indistinguishable once each glyph is isolated
and resized, a substantial fraction of the 62-way decisions cannot be made
from a $32\times32$ image at all. This shows in the per-class results:
uppercase letters, which have the fewest look-alikes, average 0.534 F1, while
digits (0.414) and lowercase letters (0.415) trail them, and the worst
individual classes are precisely the ambiguous ones (\texttt{x},
\texttt{I}, \texttt{k}, \texttt{0}, \texttt{z}, \texttt{s}). Second, the
staged tuning improved cross-validated accuracy by 4.5 points but did
\emph{not} improve held-out test accuracy over the untuned baseline
(46.48\% vs 48.83\%), and the tuned model memorises its training set almost
perfectly. The limiting factor at this point is data quantity and
representation, not hyperparameter choice -- which is why augmentation or a
convolutional model, rather than further tuning of a dense network, is the
right next step.

\section{References}
\begin{enumerate}[nosep]
\item F. Rosenblatt, ``The perceptron: A probabilistic model for information
storage and organization in the brain,'' \emph{Psychological Review}, vol. 65,
no. 6, pp. 386--408, 1958.
\item D. E. Rumelhart, G. E. Hinton, and R. J. Williams, ``Learning
representations by back-propagating errors,'' \emph{Nature}, vol. 323,
pp. 533--536, 1986.
\item D. P. Kingma and J. Ba, ``Adam: A Method for Stochastic Optimization,''
\emph{3rd International Conference on Learning Representations (ICLR)}, 2015.
\item Scikit-learn Developers, ``Neural network models (supervised),''
scikit-learn documentation. [Online]. Available:
\url{https://scikit-learn.org/stable/modules/neural_networks_supervised.html}
\item D. Dave, ``English Handwritten Characters Dataset,'' Kaggle. [Online].
Available: \url{https://www.kaggle.com/datasets/dhruvildave/english-handwritten-characters-dataset}
\end{enumerate}

\end{document}
